\chapter{SYSTEM ANALYSIS}

\section{Module Description}

The Medora Hospital Management System is divided into multiple functional modules, each responsible for managing specific hospital operations efficiently and securely. The major modules are described below:

\begin{itemize}

\item \textbf{User Management Module:}
This module manages registration, authentication, and role-based access control for all users including Admin, Receptionist, Doctor, Patient, Pharmacist, and Lab Staff. It ensures secure login and restricts access to system features based on user roles.

\item \textbf{Reception Module:}
This module handles patient registration, appointment booking, token generation, and consultation billing. Receptionists can schedule appointments based on doctor availability and manage patient details.

\item \textbf{Doctor Module:}
This module enables doctors to view appointments, access patient history, enter diagnoses, generate digital prescriptions, and request laboratory tests. It ensures accurate documentation of consultations.

\item \textbf{Pharmacy Module:}
The pharmacy module manages medicine inventory, prescription verification, billing, and dispensing. It also supports direct medicine purchase where patients can upload external prescriptions or purchase medicines without consultation.

\item \textbf{Laboratory Module:}
This module handles lab test requests, test result uploads, and report generation. Lab staff update test status and maintain laboratory records securely.

\item \textbf{Billing Module:}
This module generates invoices for consultations, lab tests, and medicine purchases. It calculates doctor fees, additional charges, and tracks payment status.

\item \textbf{Appointment Management Module:}
This module manages doctor schedules, appointment slots, token numbers, and appointment status updates.

\end{itemize}

\section{Business Rules}

Business rules define operational constraints to ensure system consistency and reliability.

\textbf{User and Access Rules}
\begin{itemize}
\item Every user must register before accessing the system.
\item Role-based access control restricts access to authorized features only.
\item Only Admin can manage staff accounts.
\end{itemize}

\textbf{Appointment Rules}
\begin{itemize}
\item A patient cannot book multiple appointments for the same doctor at the same time.
\item Appointment booking is based on doctor availability.
\item Consultation payment must be completed before marking appointment as finished.
\end{itemize}

\textbf{Prescription Rules}
\begin{itemize}
\item Only doctors can generate prescriptions.
\item Prescriptions are linked to specific patients.
\item Pharmacists can dispense medicines only for valid prescriptions.
\end{itemize}

\textbf{Pharmacy Rules}
\begin{itemize}
\item Medicine stock must be verified before dispensing.
\item Patients can purchase medicines directly without doctor consultation.
\item Uploaded prescriptions must be validated by the pharmacist.
\end{itemize}

\textbf{Laboratory Rules}
\begin{itemize}
\item Lab tests must be linked to a patient or appointment.
\item Only lab staff can upload results.
\item Finalized reports cannot be modified.
\end{itemize}

\section{Analysis of Architecture}

The Medora Hospital Management System follows the Laravel MVC (Model–View–Controller) architecture to ensure structured and maintainable development.

\begin{enumerate}

\item \textbf{Model Layer:}
Handles database interactions using Eloquent ORM. It defines relationships between entities such as User, Doctor, Patient, Appointment, Prescription, LabTest, Billing, and Medicine.

\item \textbf{View Layer:}
Developed using Blade templating engine to display role-specific dashboards and dynamic data.

\item \textbf{Controller Layer:}
Manages request handling, validation, business logic processing, and user redirection.

\item \textbf{Database Layer:}
Uses MySQL for storing structured hospital data including appointments, prescriptions, lab reports, and billing.

\item \textbf{Security Features:}
Implements password hashing, CSRF protection, middleware-based role gating, and soft deletes for data safety.

\end{enumerate}

\section{Feasibility Analysis}

\subsection{Technical Feasibility}

The system is developed using Laravel framework and MySQL database, which are reliable, scalable, and widely supported technologies. The application runs on Apache or Nginx servers and can be deployed easily in a hospital environment. The development tools ensure maintainability and future scalability.

\subsection{Economic Feasibility}

The project is cost-effective as it uses open-source technologies such as Laravel and MySQL. Development and maintenance costs are minimal, and automation reduces manual paperwork and staffing overhead.

\subsection{Operational Feasibility}

The system is user-friendly and aligns with hospital workflows. Receptionists, doctors, pharmacists, and lab staff can easily adapt to the digital system. The structured modules ensure smooth daily operations.

\section{System Environment}

\subsection{Software Environment}

\textbf{Operating System:} Windows 10/11 or Linux \\
\textbf{Backend Framework:} Laravel 12 \\
\textbf{Programming Language:} PHP 8.2+ \\
\textbf{Frontend:} Blade, HTML, CSS, JavaScript \\
\textbf{Database:} MySQL 8.0+ \\
\textbf{Web Server:} Apache / Nginx \\
\textbf{Development Tools:} VS Code, Composer, Git \\

\subsection{Hardware Environment}

\textbf{Processor:} Intel i5 12th Gen \\
\textbf{Memory (RAM):} 16 GB \\
\textbf{Storage:} 512 GB SSD \\
\textbf{Operating System:} Windows 11 \\

\section{Actors and Roles}

\textbf{Admin}
\begin{itemize}
\item Manage hospital staff accounts.
\item Monitor appointments and system reports.
\item Handle complaints and configurations.
\end{itemize}

\textbf{Receptionist}
\begin{itemize}
\item Register patients.
\item Book appointments.
\item Generate tokens.
\item Process billing.
\end{itemize}

\textbf{Doctor}
\begin{itemize}
\item View appointments.
\item Diagnose patients.
\item Generate prescriptions.
\item Request lab tests.
\end{itemize}

\textbf{Pharmacist}
\begin{itemize}
\item Verify prescriptions.
\item Dispense medicines.
\item Manage inventory.
\item Handle direct medicine purchases.
\end{itemize}

\textbf{Lab Staff}
\begin{itemize}
\item View lab requests.
\item Upload test results.
\item Update lab status.
\end{itemize}

\textbf{Patient}
\begin{itemize}
\item Register and login.
\item Book appointments.
\item View prescriptions and lab reports.
\item Upload external prescriptions.
\item Purchase medicines directly.
\end{itemize}